\begin{initiativkarte}{MakerSpace – Fakultät EMI (HAW Hamburg)}

  % Bild (optional)
  \IfFileExists{bilder/makerspace.jpg}{%
    \begin{center}
      \includegraphics[width=\linewidth,height=5cm,keepaspectratio]{bilder/makerspace.jpg}
    \end{center}
    \vspace{0.4cm}
  }{}

  % Kurzbeschreibung
  \textit{\textcolor{hawgray}{Der MakerSpace der Fakultät EMI ist ein offener Lern- und Arbeitsraum, in dem Studierende praktische Technik-, Elektronik- und Hardwareprojekte umsetzen können.}}

  \vspace{0.5cm}
  \hawrule

  % Beschreibung
  \textbf{\textcolor{hawblue}{Was machen wir?}}\\[0.3cm]

  Der MakerSpace an der Fakultät Elektro-, Medien- und Informationstechnik (EMI) der HAW Hamburg bietet Studierenden die Möglichkeit, eigene technische Projekte praktisch umzusetzen.

  Im Mittelpunkt steht das „Learning by Doing“ mit moderner Labor- und Entwicklungsausstattung.

  Zu den typischen Aktivitäten gehören:
  \begin{itemize}
      \item Entwicklung und Aufbau elektronischer Schaltungen
      \item Löten und Arbeiten mit Hardware
      \item Programmierung von Mikrocontrollern (z. B. Arduino, ESP32)
      \item Nutzung von Oszilloskopen, Funktionsgeneratoren und weiteren Laborgeräten
      \item Einsatz von 3D-Druck für Prototypen
      \item Teilnahme an technischen Workshops
  \end{itemize}

  Der MakerSpace richtet sich an Studierende, die praktische Erfahrung in Elektronik, Embedded Systems und Prototyping sammeln möchten – unabhängig vom Vorwissen.

  \vspace{0.5cm}

  % Tags
  \tag{Technik}
  \tag{Elektronik}
  \tag{Programmierung}
  \tag{Labor}
  \tag{Prototyping}

  \vspace{0.6cm}
  \hawrule

  % Infozeilen
  \infozeile{\faUsers}{Zielgruppe: Studierende der Fakultät EMI (offen für Interessierte)}
  \infozeile{\faGraduationCap}{Semester: Ab dem 1. Semester (Einsteiger-Workshops möglich)}
  \infozeile{\faMapMarker*}{Ort: HAW Hamburg, Berliner Tor 7 (EMI Labore)}
  \infozeile{\faGlobe}{Website: \href{https://www.haw-hamburg.de/elektro-medien-und-informationstechnik/labore/makerspace/}{HAW MakerSpace}}

\end{initiativkarte}